% Faithful typesetting transcription of necessary formal Section2 definitions.
\section*{2 Preliminaries}
\paragraph{Notations.} We use the following notation: $[n]=\{1,2,\ldots,n\}$. Let $S^{n\times m}$ denote the set of matrices of size $n\times m$ with elements from $S$. We denote by $0_{n\times m}$ the zero matrix of size $n\times m$ and by $1_{n\times m}$ the matrix of the same size containing only ones. For square matrices $A_1,\ldots,A_k$ we denote a diagonal block matrix with blocks $A_1,\ldots,A_k$ by $\operatorname{diag}(A_1,\ldots,A_k)$. For $x\in\{0,1,\ldots,2^k-1\}$ we denote a vector $(a_0,\ldots,a_{k-1})\in\{0,1\}^k$ such that $x=\sum_{i=0}^{k-1}a_i2^i$ by $\operatorname{bin}_k(x)$, i.e. $(a_0,\ldots,a_{k-1})$ is the reversed binary representation of $x$. For vectors $v_1,\ldots,v_n$ from a vector space over a field $\mathbb F$ we denote their linear span by $\operatorname{Span}(v_1,\ldots,v_n)$. We use coordinate-wise comparison of strings from $\{0,1\}^n$, i.e. for $x,y\in\{0,1\}^n$ we write $x\le y$ iff $x_i\le y_i$ for each $i\in[n]$. We denote the set of variables of a CNF-formula $\phi$ by $\operatorname{Vars}(\phi)$.
\paragraph{Communication complexity.} We briefly recall some notions of communication complexity. For formal definition and details we refer to [24].
In the classic two-party randomized communication protocol with public randomness, Alice and Bob cooperate to compute a relation $Q\subseteq X\times Y\times Z$: Alice has an input $x\in X$ and Bob has an input $y\in Y$, their goal is to compute $z\in Z$ such that $(x,y,z)\in Q$. We assume that Alice and Bob have access to an arbitrary large random string of bits that is common for Alice and Bob. Let for every $x\in X$ and $y\in Y$, $R^\delta_{\mathrm{pub}}(Q,x,y)$ denote the minimal number of bits Alice and Bob need to transmit between each other so they both find a $z\in Z$ such that $(x,y,z)\in Q$ with probability at least $1-\delta$ taken over the values of the common random string. And $R^\delta_{\mathrm{pub}}(Q):=\max_{x\in X,y\in Y}R^\delta_{\mathrm{pub}}(Q,x,y)$.
We also consider multiparty communication protocols in the number on forehead (NOF) model that extends two-party protocols for an arbitrary number of parties. In this setting $k$ parties cooperate to compute a relation $Q\subseteq X_1\times X_2\times\cdots\times X_k\times Y$. The $i$th party has $x_i\in X_i$ written on their forehead so they know all $x_j$ for $j\ne i$, their goal is to compute $y\in Y$ such that $(x_1,x_2,\ldots,x_k,y)\in Q$. The parties communicate by taking turns broadcasting messages to all other parties until all parties learn the value of $y\in Y$ such that $(x_1,\ldots,x_k,y)\in Q$. In this model we also assume that all parties have access to a common random string of bits. Let $R^\delta_{\mathrm{pub}}(Q,x_1,\ldots,x_k)$ for $x_1\in X_1,\ldots,x_k\in X_k$ denote the minimal total number of bits transmitted until each party learns $y\in Y$ such that $(x_1,\ldots,x_k,y)\in Q$ with probability at least $1-\delta$ taken over the set of values of the random string of bits. Also, let $R^\delta_{\mathrm{pub}}(Q):=\max_{x_1\in X_1,\ldots,x_k\in X_k}R^\delta_{\mathrm{pub}}(Q,x_1,\ldots,x_k)$.
Let $f$ be a function from $X_1\times X_2\times\cdots\times X_k\to Y$. Then $R^\delta_{\mathrm{pub}}(f)$ denotes $R^\delta_{\mathrm{pub}}(Q_f)$, where $Q_f=\{(x_1,x_2,\ldots,x_k,y)\mid f(x_1,\ldots,x_k)=y\}$.
We prove communication complexity lower bounds by reduction from different versions of the set disjointness problem. $\mathrm{DISJ}_{k,n}$ is a function $\{0,1\}^{kn}\to\{0,1\}$ such that for every $x_1,\ldots,x_k\in\{0,1\}^n$ the following holds:
\[\mathrm{DISJ}_{k,n}(x_1,\ldots,x_k)=\bigwedge_{j=1}^n\neg\left(\bigwedge_{i=1}^k(x_i)_j\right).\]
Let us define the communication promise problem $\mathrm{UDISJ}_{k,n}$ in the $k$-party NOF model. For each $i\in[k]$ the string $x_i$ is written on the forehead of the $i$th party, it is guaranteed that there exists at most one index $j\in[n]$ such that for every $i\in[k]$, $(x_i)_j=1$. The goal is to compute $\mathrm{DISJ}_{k,n}(x_1,\ldots,x_k)$.
\begin{authorclaim}{Theorem7 ([31,32])}
$R^{1/3}_{\mathrm{pub}}(\mathrm{UDISJ}_{k,n})=\Omega(\sqrt n/(2^kk))$.
\end{authorclaim}
For $k=2$ we omit the first index: $\mathrm{DISJ}_n=\mathrm{DISJ}_{2,n}$; in this case Theorem7 may be improved.
\begin{authorclaim}{Theorem8 ([16])}
$R^{1/3}_{\mathrm{pub}}(\mathrm{DISJ}_n)\ge R^{1/3}_{\mathrm{pub}}(\mathrm{UDISJ}_{2,n})=\Omega(n)$.
\end{authorclaim}
\paragraph{Proof complexity.} We consider refutational proof systems for the language of unsatisfiable CNF-formulas UNSAT. A refutation of $\phi\in\mathrm{UNSAT}$ in a proof system $\Pi$ is a sequence of Boolean functions (proof lines) such that each proof line either represents a clause of $\phi$ or derived from previous proof lines in the sequence via some sound inference rules. The last line of the proof is identically zero function. A proof system $\Pi$ is defined by a representation of proof lines and by a set of admissible inference rules. It is required that the inference rules are polynomially verifiable i.e. there exists an algorithm that checks whether it is legitimate to derive a line $L_0$ from the lines $L_1,\ldots,L_k$.
For example, in the Resolution proof lines are represented by clauses and the only inference rule is the resolution rule that allows deriving a clause $A\vee B$ from the clauses $A\vee x$ and $A\vee\neg x$.
The size of a proof is the total size of all representations of proof lines in the proof. The length of a proof is the number of proof lines in it.
A tree-like proof is such a proof that every its line can be used as a premise of a rule at most once. For each proof system, we can also consider its tree-like version where all proofs are constrained to be tree-like.
We also consider semantic refutational proof systems, where we drop the requirement for polynomial verification of inference rules i.e. we allow to derive any sound consequence from the premises. For such systems it is crucial to bound fan-in i.e. the number of the premises from which each proof line can be derived, otherwise, it would be possible to derive a contradiction from the clauses of the initial formula immediately. For example, it is well-known that Resolution is polynomially equivalent to a semantic proof system with fan-in 2 operating with clauses.
A lower bound on the proof size in a semantic proof system implies a lower bound on the proof size in its syntactic counterpart because a syntactic proof is always a semantic proof that operates with the same class of proof lines.
We define semantic $\mathrm{Res}(\oplus)$ as a semantic proof system with fan-in 2 that operates with linear clauses. A linear clause is a disjunction of linear equations over $\mathbb F_2$: $\bigvee_{i=1}^k(f_i=a_i)$, where $f_i$ is a linear form over $\mathbb F_2$ and $a_i\in\mathbb F_2$. Notice that an ordinary clause $\bigvee_{i\in P}x_i\vee\bigvee_{j\in N}\neg x_j$ can be represented by the linear clause $\bigvee_{i\in P}(x_i=1)\vee\bigvee_{j\in N}(x_j=0)$. For definition of syntactic version of $\mathrm{Res}(\oplus)$ we refer to [15]; it is also proved there that syntactic and semantic $\mathrm{Res}(\oplus)$ are polynomially equivalent.
We define semantic $\mathrm{Res}(\mathrm{PC}_d)$ as a semantic proof system with fan-in 2 that operates with disjunctions of equations of type $f=0$, where $f$ is a degree-$d$ polynomial over $\mathbb F_2$. Notice that semantic $\mathrm{Res}(\mathrm{PC}_1)$ is exactly semantic $\mathrm{Res}(\oplus)$. For the definition of the syntactic version of $\mathrm{Res}(\mathrm{PC}_d)$ we refer to [19].
Following [1] we define $\mathrm{Th}(k)$ as a semantic proof system with fan-in 2 that operates with polynomial inequalities $g\ge0$, where $g$ is a polynomial of degree at most $k$ with integer coefficients and Boolean variables. A clause $\bigvee_{i\in P}x_i\vee\bigvee_{j\in N}\neg x_j$ can be represented by an inequality $\sum_{i\in P}x_i+\sum_{j\in N}(1-x_j)-1\ge0$.
\paragraph{Proof complexity and communication complexity.} For an unsatisfiable CNF-formula $\phi$ we define the communication problem $\operatorname{Search}(\phi)$. Search$(\phi)$ is the following problem: given an assignment of the variables of the unsatisfiable CNF $\phi$, find a clause that is falsified by this assignment. It is assumed that variables of $\phi$ are somehow partitioned between the parties.
Following the paper [9] we consider a semantic proof system $\mathrm T^{\mathrm{cc}}(k,c)$ that models many interesting syntactic and semantic proof systems. The proof lines in $\mathrm T^{\mathrm{cc}}(k,c)$ can be arbitrary Boolean functions having the following property: for every proof line $C$ and every partition of variables of $C$ between $k$ parties, the NOF $k$-party randomized communication complexity of $C$ is at most $c$ w.r.t. this partition. We also define a semantic proof system $\mathrm T^{\mathrm{cc}}_{\mathrm{os}}(k,c)$ that is a subsystem of $\mathrm T^{\mathrm{cc}}(k,c)$ with the restriction that a communication protocol for proof lines must have a one-sided error: if the value of a proof line is zero, then the protocol should return zero with probability 1.
For example, $\mathrm T^{\mathrm{cc}}(2,2)$ simulates Resolution; $\mathrm T^{\mathrm{cc}}(2,O(1))$ simulates $\mathrm{Res}(\oplus)$ [22,15]; $\mathrm T^{\mathrm{cc}}(k,O(k^3\log^2n))$, where $n$ is the number of variables in a refuted formula, simulates $\mathrm{Th}(k-1)$ [9]. In Section3 we show that $\mathrm T^{\mathrm{cc}}_{\mathrm{os}}(d+1,O(1))$ simulates $\mathrm{Res}(\mathrm{PC}_d)$.
The following connection between the communication complexity of Search$(\phi)$ and tree-like proof complexity of $\phi$ is known.
\begin{authorclaim}{Lemma9 ([1,9])}
If a CNF formula $\phi$ has a tree-like $\mathrm T^{\mathrm{cc}}(k,c)$ refutation of length $\ell$ then, over any $k$-partition of the variables, there is a randomized bounded-error $k$-party NOF protocol for Search$(\phi)$ with communication cost $O(c\cdot\log\ell\log\log\ell)$.
\end{authorclaim}
In Section3 we show that for $\mathrm T^{\mathrm{cc}}_{\mathrm{os}}(k,c)$ the bound can be improved, see Remark14.
\paragraph{Basic formulas.} A CNF formula $\mathrm{PHP}_n^m$ encodes the pigeonhole principle; $\mathrm{PHP}_n^m$ states that it is possible to put $m$ pigeons into $n$ holes such that every pigeon flies to at least one hole and at most one pigeon flies to each hole. $\mathrm{PHP}_n^m$ depends on variables $p_{i,j}$ for $i\in[m]$ and $j\in[n]$ and $p_{i,j}=1$ iff the $i$-th pigeon flies to the $j$-th hole. $\mathrm{PHP}_n^m$ is the conjunction of $m(m-1)n/2$ hole axioms and $m$ pigeons axioms. For every $i\in[m]$ $\mathrm{PHP}_n^m$ contains a pigeon axiom $(p_{i,1}\vee p_{i,1}\vee\cdots\vee p_{i,n})$. And for every $j\in[n]$ and every $k\ne\ell\in[n]$, $\mathrm{PHP}_n^m$ contains a hole axiom $(\neg p_{k,j}\vee\neg p_{\ell,j})$. $\mathrm{PHP}_n^m$ is unsatisfiable iff $m>n$.
For an undirected graph $G(V,E)$, the formula $\mathrm{PM}_G$ encodes in CNF that $G$ has a perfect matching. The formula $\mathrm{PM}_G$ has $|E|$ variables, each of them corresponds to an edge of $G$, $x_e$ is the variable corresponding to $e\in E$.
\[\mathrm{PM}_G=\bigwedge_{v\in V}\left(\left(\bigvee_{e\text{ is incident to }v}x_e\right)\wedge\bigwedge_{e_1\ne e_2\text{ are incident to }v}(\neg x_{e_1}\vee\neg x_{e_2})\right).\]
$\mathrm{PM}_G$ is unsatisfiable iff $G$ does not have a perfect matching.
\begin{authorclaim}{Theorem10 ([26])}
Let $G$ be a graph with $n$ vertices, which has no perfect matching. Then the formula $\mathrm{PM}_G$ has a tree-like $\mathrm{Res}(\oplus)$ refutation of size $2^{O(n)}$.
\end{authorclaim}
\begin{authorclaim}{Proposition11 ([14])}
Let $G$ be a graph with an odd number of vertices. Then the formula $\mathrm{PM}_G$ has a tree-like $\mathrm{Res}(\oplus)$ refutation of size $\operatorname{poly}(n)$.
\end{authorclaim}
The binary pigeonhole principle $\mathrm{BPHP}_{2^\ell}^m$ states that there are $m$ different $\ell$-bit binary strings $s_1,s_2,\ldots,s_m$. $\mathrm{BPHP}_{2^\ell}^m$ has $m\ell$ variables corresponding to the bits of $s_i$ for $i\in[m]$. Then $\mathrm{BPHP}_{2^\ell}^m=\bigwedge_{i\ne j\in[m]}s_i\ne s_j$, where the predicate $s_i\ne s_j$ is encoded as a $2\ell$-CNF formula of size $2^\ell$ as follows: $\bigwedge_{\alpha\in\{0,1\}^\ell}(s_i\ne\alpha\vee s_j\ne\alpha)$; notice that the predicate $(s_i\ne\alpha\vee s_j\ne\alpha)$ can be represented by a clause with $2\ell$ literals. If $m>2^\ell$, then the formula $\mathrm{BPHP}_{2^\ell}^m$ is unsatisfiable.
Let $\phi$ be a CNF formula with $n$ variables, and $g:\{0,1\}^k\to\{0,1\}$ be a Boolean function. Then $\phi\circ g$ denotes a CNF formula on $kn$ variables that represents $\phi(g(\vec x_1),g(\vec x_2),\ldots,g(\vec x_n))$, where $\vec x_i$ denotes a vector of $k$ new variables. $\phi\circ g$ is constructed by applying the substitution to every clause $C$ of $\phi$ and converting the resulting function $C\circ g$ to CNF in some fixed way.
